\documentclass[11pt]{article}

\usepackage{amsmath,amssymb}
\usepackage{algorithm}
\usepackage[noend]{algpseudocode}
\usepackage{geometry}
\usepackage{xurl}
\usepackage{booktabs}
\geometry{margin=1in}

\title{Construction of SPF Datasets for the ME3AI Project}
\date{}

\begin{document}
\maketitle

\begin{abstract}
This document specifies the procedure for constructing Survey of Professional
Forecasters (SPF) datasets used in this project. The SPF is a quarterly survey
of macroeconomic forecasts conducted by the Federal Reserve Bank of
Philadelphia; data are provided as Excel workbooks (median and mean levels and
growth rates, probabilities, dispersion, and microdata). We describe how to
obtain the official releases, which variables and vintages to use, how to
validate and normalize the data into the project's cleaned third normal form
(3NF) structure, and how to link them to the pipeline input and output folders.
This document serves as the single source of truth for SPF dataset construction
and reproducibility.
\end{abstract}

\begin{table}[h]
\centering
\caption{Notation and terms used in this document.}
\label{tab:notation}
\begin{tabular}{@{}ll@{}}
\toprule
Symbol or term & Definition \\
\midrule
SPF & Survey of Professional Forecasters \\
3NF & Third normal form (database normalization) \\
BEA & Bureau of Economic Analysis \\
NIPA & National income and product accounts \\
     & (BEA advance report) \\
\texttt{var\_path} & URL path segment for a variable (lowercase; \\
     & e.g., \texttt{ngdp}, \texttt{cpi10}) \\
VAR & Placeholder for variable code in filenames (e.g., NGDP, CPI10) \\
\texttt{overwrite} & Policy: \texttt{overwrite} or \texttt{skip-if-exists} \\
\bottomrule
\end{tabular}
\end{table}

\section{Data source and scope}

The Survey of Professional Forecasters (SPF) is a quarterly survey of
professional forecasters, formerly run by the American Statistical Association
and the National Bureau of Economic Research (from 1968:Q4) and taken over by
the Federal Reserve Bank of Philadelphia in 1990:Q2. The Philadelphia Fed
publishes documentation and data files; the canonical reference for variable
definitions, file layout, and timing is the \emph{SPF Documentation}
(Philadelphia Fed, latest update March 2024), held in this project as
\texttt{reference/spf-documentation.pdf}.

The survey is timed to the release of the Bureau of Economic Analysis (BEA)
advance report of the national income and product accounts (NIPA):
questionnaires are sent after that report; deadlines fall in the second to
third week of the middle month of each quarter; results are released before
the BEA second report. For the construction of our datasets, we use the
survey date (year and quarter) as the primary time index. The first survey
fully conducted in real time by the Philadelphia Fed is 1990:Q3. Earlier
vintages (1968:Q4 to 1990:Q2) may be used with the caveats noted in the
official documentation.

\section{Published files and variables}

The Philadelphia Fed provides SPF data \emph{by variable}: each variable has a
dedicated data page from which Excel (and other) files can be downloaded. The
variable index is at \url{https://www.philadelphiafed.org/surveys-and-data/data-files};
clicking a variable (e.g., NGDP, CPI10, RGDP, UNEMP) leads to a page at
\[
\texttt{https://www.philadelphiafed.org/surveys-and-data/}\textit{var\_path},
\]
where $\textit{var\_path}$ is the variable identifier in lowercase (e.g.,
\texttt{ngdp}, \texttt{cpi10}, \texttt{rgdp}). The sole exception is CPI
inflation, which uses \texttt{cpi-spf}. Variable names and paths are listed on
the data-files index (e.g., ``10-Year CPI Inflation Rate (CPI10)'' maps to
\texttt{/cpi10}).

On each variable page, the following file types are typically offered
(availability may vary by variable):

\begin{itemize}
  \item \textbf{Documentation}: \\
    Link to the SPF documentation PDF (same for all variables).
  \item \textbf{Measures of Cross-Sectional Forecast Dispersion}: \\
    \texttt{Dispersion\_}\textit{VAR}\texttt{.xlsx} (e.g.,
    \texttt{Dispersion\_CPI10.xlsx}, \texttt{Dispersion\_NGDP.xlsx}).
  \item \textbf{Median Responses}: \\
    \texttt{Median\_}\textit{VAR}\texttt{\_Level.xlsx} (e.g.,
    \texttt{Median\_NGDP\_Level.xlsx}, \texttt{Median\_CPI10\_Level.xlsx}).
  \item \textbf{Mean Responses}: \\
    \texttt{Mean\_}\textit{VAR}\texttt{\_Level.xlsx}.
  \item \textbf{Annualized Percent Change of Median/Mean Responses}: \\
    \texttt{Median\_}\textit{VAR}\texttt{\_Growth.xlsx} and
    \texttt{Mean\_}\textit{VAR}\texttt{\_Growth.xlsx}. For variables that have
    growth rates (e.g., NGDP, RGDP); not all variables such as CPI10 offer
    growth files.
  \item \textbf{Individual Responses}: \\
    \texttt{Individual\_}\textit{VAR}\texttt{.xlsx}.
\end{itemize}

Filenames use the variable code as on the website (e.g., NGDP, CPI10). The
automation must fetch the variable page, extract the download links for the
desired file types, and save each file to \texttt{data-pipeline/input/} under
that same filename (no timestamps). Variable names, horizons, and caveats are
defined in the SPF documentation; we record the exact variable names and
definitions in a project codebook.

\section{Download automation: inputs and outputs}

We automate retrieval of SPF files from the Philadelphia Fed website so that
raw files are saved to \texttt{data-pipeline/input/} without manual download.
Given a variable name (e.g., CPI10, NGDP), the data are hosted on that
variable's page at \url{https://www.philadelphiafed.org/surveys-and-data/}%
$\textit{var\_path}$.
The automation fetches the page, parses the HTML for the desired download
links (Dispersion, Median Level, Mean Level, Median Growth, Mean Growth,
Individual), and downloads each file to the input folder under the same
filename as on the site (e.g., \texttt{Median\_CPI10\_Level.xlsx}). When
parsing links, the implementation must match \texttt{<a>} tags where
\texttt{href} may appear after other attributes (e.g.\ \texttt{class}), since
the Philadelphia Fed markup may use that order.

\paragraph{Inputs.}
A list of variable names (e.g., CPI10, RGDP, UNEMP, NGDP) as on the
Philadelphia Fed data-files index. Optional inputs: which file types to
download (e.g., median level, mean level, dispersion, individual); overwrite
policy $\texttt{overwrite} \in \{\texttt{overwrite}, \texttt{skip-if-exists}\}$
(whether to overwrite an existing file or skip); output directory (default
\texttt{data-pipeline/input/}).

\paragraph{Outputs.}
Each selected file is saved under \texttt{data-pipeline/input/} with the same
filename as on the website (no timestamps or version suffixes). Each file is
downloaded at most once per variable; if multiple variables were to share the
same filename, the implementation should deduplicate so that the file is
written once.

Three logical functions suffice: (1) given a variable name, return the variable
page URL (lowercase path, with exception CPI $\to$ \texttt{cpi-spf}); (2) given
a variable page URL and requested file types, fetch the page and return a list
of (filename, download URL) pairs; (3) given a download URL and local path,
download the file subject to the overwrite policy. The top-level procedure
loops over the requested variables, obtains the link list for each, and
downloads each file to \texttt{data-pipeline/input/}.

\section{Pipeline and 3NF}

Our data pipeline (see \texttt{docs/workflow/data\_pipeline.md}) has three
stages: \textbf{input} (raw data), \textbf{cleaned} (third normal form, 3NF),
and \textbf{output} (analysis results). SPF data enter as raw files in the
input folder and are transformed into cleaned tables that satisfy 3NF.

\paragraph{Input.}
Place the official Excel workbooks in the pipeline input folder. Do not
modify the original files; preserve them as the source of truth.

\paragraph{Cleaned.}
The current implementation focuses on \textbf{individual (microdata)} forecasts
only. Cleaned output is CSV: \texttt{forecast\_individual.csv} in \textbf{wide}
form (primary key: survey\_year, survey\_quarter, variable, forecaster\_id; one
column per horizon, all horizons in the same row) and \texttt{forecaster\_survey.csv}
(primary key: survey\_year, survey\_quarter, forecaster\_id; \texttt{industry}).
No separate survey table is used; the set of (survey\_year, survey\_quarter) is
derivable from \texttt{forecast\_individual}. Level/mean/median and probability
tables may be added in a later stage.

\paragraph{Output.}
Analyses read from the cleaned folder and write results to the output
folder, organized by analysis type.

Before treating a constructed SPF dataset as final: verify that raw files match
the expected Philadelphia Fed release; confirm variable names and column
layouts align with the SPF documentation; after transformation, check row
counts and primary keys; record data source, retrieval date, and any
filtering in the codebook or a README in the cleaned folder.

\section{Algorithms: download automation}

\begin{algorithm}[H]
\caption{\textsc{VariablePageURL}$(\textit{var})$}
\begin{algorithmic}[1]
\State \textbf{Require} $\textit{var}$: variable name as on Philadelphia Fed data-files index (e.g., NGDP, CPI10).
\State $\textit{path} \gets \text{lowercase}(\textit{var})$.
\If{$\textit{var} = \texttt{CPI}$}
    \State $\textit{path} \gets \texttt{cpi-spf}$. \Comment{exception per Philadelphia Fed URL}
\EndIf
\State \Return \texttt{https://www.philadelphiafed.org/surveys-and-data/} $\oplus$ $\textit{path}$.
\end{algorithmic}
\end{algorithm}

\begin{algorithm}[H]
\caption{\textsc{GetDownloadLinks}$(\textit{page\_url}, \textit{file\_types})$}
\begin{algorithmic}[1]
\State \textbf{Require} $\textit{file\_types}$: subset of
\Statex \quad $\{\texttt{dispersion}, \texttt{median\_level}, \texttt{mean\_level}, \texttt{median\_growth}, \texttt{mean\_growth}, \texttt{individual}\}$, and optionally \texttt{documentation}.
\State Fetch HTML at $\textit{page\_url}$.
\State Parse page for download links: match \texttt{<a>} tags allowing
\Statex \quad \texttt{href} to appear after other attributes (e.g.\ \texttt{class}); match link text or href to each requested type (e.g., ``Median Responses'' $\to$ \texttt{Median\_VAR\_Level.xlsx}, ``Measures of Cross-Sectional Forecast Dispersion'' $\to$ \texttt{Dispersion\_VAR.xlsx}).
\State Build list $L$ of $(\textit{filename}, \textit{download\_url})$ for each matched link; extract $\textit{filename}$ from the URL or link (stable name, e.g., \texttt{Median\_NGDP\_Level.xlsx}).
\State \Return $L$.
\end{algorithmic}
\end{algorithm}

\begin{algorithm}[H]
\caption{\textsc{DownloadOneFile}$(\textit{download\_url}, \textit{outpath}, \texttt{overwrite})$}
\begin{algorithmic}[1]
\State \textbf{Require} $\texttt{overwrite} \in \{\texttt{overwrite}, \texttt{skip-if-exists}\}$.
\If{$\texttt{overwrite} = \texttt{skip-if-exists}$ and file exists at $\textit{outpath}$}
    \State \Return (skip).
\EndIf
\State Download from $\textit{download\_url}$ to $\textit{outpath}$ (overwrite if file exists when $\texttt{overwrite} = \texttt{overwrite}$).
\State \Return $\textit{outpath}$.
\end{algorithmic}
\end{algorithm}

\begin{algorithm}[H]
\caption{\textsc{DownloadByVariableNames}$(\textit{variable\_names}, \textit{file\_types}, \textit{out\_dir}, \texttt{overwrite})$}
\begin{algorithmic}[1]
\State \textbf{Require} $\textit{out\_dir} = \texttt{data-pipeline/input/}$ (fixed by project layout).
\State \textbf{Require} $\textit{file\_types}$: which file types to download per variable (e.g., median level, mean level, dispersion, individual).
\State $\textit{written} \gets \emptyset$. \Comment{optional: track filenames already written to avoid duplicate writes}
\For{each $v$ in $\textit{variable\_names}$}
    \State $\textit{page\_url} \gets \textsc{VariablePageURL}(v)$.
    \State $L \gets \textsc{GetDownloadLinks}(\textit{page\_url}, \textit{file\_types})$.
    \For{each $(\textit{filename}, \textit{url})$ in $L$}
        \State $\textit{outpath} \gets \textit{out\_dir}$ joined with $\textit{filename}$.
        \State \textsc{DownloadOneFile}$(\textit{url}, \textit{outpath}, \texttt{overwrite})$.
        \State Add $\textit{outpath}$ to $\textit{written}$ (if not already present).
    \EndFor
\EndFor
\State \Return $\textit{written}$ (list of paths in \texttt{data-pipeline/input/}).
\end{algorithmic}
\end{algorithm}

\section{Algorithm: construct cleaned datasets (individual only)}

The following procedure summarizes the construction of SPF \emph{individual}
forecast datasets from raw Excel files to cleaned CSV. Raw files
\texttt{Individual\_VAR.xlsx} are assumed present in \texttt{data-pipeline/input/}
(e.g., produced by \textsc{DownloadByVariableNames} with \texttt{individual} in
\texttt{file\_types}). Output is CSV; \texttt{forecast\_individual} is written in
wide form (one row per survey/variable/forecaster, one column per horizon).
Helper functions used by the main procedure are specified below.

\begin{algorithm}[H]
\caption{\textsc{VariableFromFilename}$(\textit{path})$}
\begin{algorithmic}[1]
\State \textbf{Require} $\textit{path}$: filename or path of form \texttt{Individual\_VAR.xlsx}.
\State $\textit{stem} \gets \text{filename without extension of }\textit{path}$.
\If{$\textit{stem}$ starts with \texttt{Individual\_}}
    \State \Return $\textit{stem}$ with prefix \texttt{Individual\_} removed (e.g., VAR).
\EndIf
\State \Return $\textit{stem}$.
\end{algorithmic}
\end{algorithm}

\begin{algorithm}[H]
\caption{\textsc{LoadIndividualSheet}$(\textit{path})$}
\begin{algorithmic}[1]
\State \textbf{Require} $\textit{path}$: path to one \texttt{Individual\_VAR.xlsx} file.
\State Load first sheet: row 1 = header; first four columns YEAR, QUARTER, ID, INDUSTRY; remaining columns = horizon names.
\State Validate: at least YEAR, QUARTER, ID present; at least one horizon column.
\State Melt to long form: id\_vars = (YEAR, QUARTER, ID, INDUSTRY), value\_vars = horizon columns; create columns \texttt{horizon}, \texttt{value}.
\State Rename YEAR $\to$ survey\_year, QUARTER $\to$ survey\_quarter, ID $\to$ forecaster\_id.
\State $\textit{variable} \gets \textsc{VariableFromFilename}(\textit{path})$; add column \texttt{variable}.
\State \Return long-form table (survey\_year, survey\_quarter, forecaster\_id, industry if present, horizon, value, variable).
\end{algorithmic}
\end{algorithm}

\begin{algorithm}[H]
\caption{\textsc{HorizonSortKey}$(\textit{col})$}
\begin{algorithmic}[1]
\State \textbf{Require} $\textit{col}$: horizon column name (e.g., CPI1, CPI10).
\State \textbf{Ensure} 10-year horizon (name ending with \texttt{10}) is ordered last.
\State \Return pair $(\textit{col}\,\text{ends with}\,\texttt{10}, \textit{col})$ for sorting (False before True, then by name).
\end{algorithmic}
\end{algorithm}

\begin{algorithm}[H]
\caption{\textsc{BuildForecastIndividual}$(\textit{df\_long})$}
\begin{algorithmic}[1]
\State \textbf{Require} $\textit{df\_long}$: concatenated long-form table from \textsc{LoadIndividualSheet}.
\State Pivot: index = (survey\_year, survey\_quarter, variable, forecaster\_id); columns = horizon; values = value; aggfunc = first.
\State Order horizon columns: sort by \textsc{HorizonSortKey} so 10-year column is last.
\State Final columns: (survey\_year, survey\_quarter, variable, forecaster\_id, horizon\_1, horizon\_2, \ldots, horizon\_10).
\State \Return wide table (key columns as integers where applicable).
\end{algorithmic}
\end{algorithm}

\begin{algorithm}[H]
\caption{\textsc{BuildForecasterSurvey}$(\textit{df\_long})$}
\begin{algorithmic}[1]
\State \textbf{Require} $\textit{df\_long}$: concatenated long-form table (must contain survey\_year, survey\_quarter, forecaster\_id, INDUSTRY or industry).
\State Select distinct (survey\_year, survey\_quarter, forecaster\_id, industry); sort by (survey\_year, survey\_quarter, forecaster\_id).
\State \Return forecaster\_survey table; primary key (survey\_year, survey\_quarter, forecaster\_id).
\end{algorithmic}
\end{algorithm}

\begin{algorithm}[H]
\caption{\textsc{CleanIndividualTo3NF}$(\textit{input\_dir}, \textit{cleaned\_dir})$}
\begin{algorithmic}[1]
\State \textbf{Require} $\textit{input\_dir}$: path to input folder; $\textit{cleaned\_dir}$: path to cleaned output folder.
\State Find all \texttt{Individual\_*.xlsx} in $\textit{input\_dir}$.
\For{each file $\textit{path}$ in that list}
    \State $\textit{df} \gets \textsc{LoadIndividualSheet}(\textit{path})$.
    \State Append $\textit{df}$ to collection of long-form tables.
\EndFor
\State $\textit{df\_long} \gets$ concatenate all long-form tables.
\State $\textit{forecast\_individual} \gets \textsc{BuildForecastIndividual}(\textit{df\_long})$.
\State $\textit{forecaster\_survey} \gets \textsc{BuildForecasterSurvey}(\textit{df\_long})$.
\State Write $\textit{forecast\_individual}$ to $\textit{cleaned\_dir}$/\texttt{forecast\_individual.csv}, $\textit{forecaster\_survey}$ to $\textit{cleaned\_dir}$/\texttt{forecaster\_survey.csv}.
\State \Return paths to cleaned CSV files.
\end{algorithmic}
\end{algorithm}

\section{Notes}

Obtaining official Excel files may be implemented by \textsc{DownloadByVariableNames}: given variable names (e.g., CPI10, NGDP), the procedure uses \textsc{VariablePageURL} to form each variable's data page URL (e.g., \url{https://www.philadelphiafed.org/surveys-and-data/cpi10}), \textsc{GetDownloadLinks} to parse the page for Individual (and optionally other) links, and \textsc{DownloadOneFile} to save each file to \texttt{data-pipeline/input/} under the same filename as on the site. The \textbf{current} cleaned schema is \texttt{forecast\_individual} and \texttt{forecaster\_survey} for individual (microdata) only; output format is CSV. Level/mean/median and probability tables are not yet in scope; table names and keys should match the pipeline and codebook when added.

\end{document}
